\documentclass[twocolumn]{aastex631}
\begin{document}
\title{The reionization ionizing-photon-budget shortfall for star-forming galaxies is not robust to systematics at z\~{}8 (optimistic systematic corner)}
\author{NebulaMind Lab (autonomous pipeline)}
\affiliation{NebulaMind Lab, \url{https://lab.nebulamind.net}}
\begin{abstract}
Reionization ionizing-photon-budget reconciliation at z\~{}8: star-forming galaxies require f\_esc=0.053 (+0.045/-0.024) to close the budget (Madau-Dickinson SFRD, log xi\_ion=25.5+/-0.15, clumping C=2-5, JWST-SFRD tail), versus indirect-proxy-inferred f\_esc=0.080 (+0.147/-0.051) from LzLCS O32/beta calibrations. Median delta(required-inferred)=-0.025 dex-frac (16-84\%: -0.168 to +0.038); 36\% of the systematic MC shows a shortfall. Conclusion: the budget CLOSES within the systematic, and the sign flips sign between the O32 and beta calibrations (calibration-driven, not robust). The result is bounded by the xi\_ion x clumping x proxy-calibration systematic, not by statistics. Generated autonomously from public data (jwst) via the NebulaMind Lab runner: a bounded, reproducible, descriptive study.
\end{abstract}
\section{Introduction}
Reionization is a pivotal period in cosmic history when the first stars and galaxies emerged, ionizing the neutral hydrogen that filled the universe. Recent studies have highlighted potential discrepancies between the observed star formation rate density (SFRD) and the required ionizing photon budget to achieve reionization [Muoz2024]. This has sparked concerns about a "photon budget crisis" and the need for alternative sources of ionizing photons, such as active galactic nuclei or X-ray binaries. To address this issue, it is essential to reconcile the cosmic SFRD with the ionizing photon budget using robust calibrations and systematic considerations [Park2022].
\section{Data and method}
In this study, we employ a literature-anchored approach to calculate the reionization-photon-budget at z\~{}8. We adopt the Madau \& Dickinson (2014) analytic fitting function for the cosmic SFRD and utilize published values for xi\_ion and O32/beta f\_esc proxy calibrations from LzLCS [Chisholm+22, Flury+22; Simmonds+24]. Our method focuses on reconciling the ionizing-photon-budget using these literature values without relying on new observational or catalog data.
\section{Result}
Our result shows that star-forming galaxies at z\~{}8 require an escape fraction of f\_esc=0.053 (+0.045/-0.024) to close the reionization photon budget, assuming a Madau-Dickinson SFRD, log xi\_ion=25.5±0.15, and clumping factor C=2-5. In contrast, indirect-proxy-inferred f\_esc from LzLCS O32/beta calibrations yields a value of 0.080 (+0.147/-0.051). The median difference between the required and inferred escape fractions is -0.025 dex (16-84\% range: -0.168 to +0.038), with 36\% of systematic Monte Carlo realizations indicating a shortfall in ionizing photons.
\begin{figure}\centering\includegraphics[width=\columnwidth]{result.png}\caption{The reionization ionizing-photon-budget shortfall for star-forming galaxies is not robust to systematics at z\~{}8 (optimistic systematic corner)}\end{figure}
\section{Caveats}
However, it is crucial to acknowledge the limitations of our approach. Our result relies on automated, single-selection, and uncalibrated measurements from published literature. The accuracy of these calibrations may be affected by factors such as sample selection biases, measurement uncertainties, and model assumptions. Furthermore, our analysis does not account for potential variations in xi\_ion or clumping factor across different galaxy populations, which could introduce additional systematic errors. Therefore, while our study provides a valuable reconciliation of the reionization photon budget, it is essential to recognize these caveats and pursue further research to refine our understanding of this critical period in cosmic history.
\section*{References}
\footnotesize
[Muoz2024] Muñoz et al. (2024). Reionization after JWST: a photon budget crisis?. bibcode 2024MNRAS.535L..37M \par [Park2022] Park et al. (2022). Calibrating excursion set reionization models to approximately conserve ionizing photons. bibcode 2022MNRAS.517..192P \par [Duncan2015] Duncan et al. (2015). Powering reionization: assessing the galaxy ionizing photon budget at z < 10. bibcode 2015MNRAS.451.2030D \par [Davies2021] Davies et al. (2021). The Predicament of Absorption-dominated Reionization: Increased Demands on Ionizing Sources. bibcode 2021ApJ...918L..35D \par [Madau2017] Madau (2017). Cosmic Reionization after Planck and before JWST: An Analytic Approach. bibcode 2017ApJ...851...50M

\end{document}
